\subsubsection*{Uniform integrality on the ray covers}
\label{subsec:theta-cover-integrality}
\hypertarget{source-theta-cover-integrality}{}

The polynomial argument requires an integral bound on the row coefficients as the squarefree ray level varies. We obtain it first on the cover, using its prescribed frames, and then descend to the level on which the determinant is evaluated.

\begin{paramdefinition}[Admitted levels and coefficient place]\label{hyp:ray-cover-integrality}
We prove the coefficient bound on the ray covers used in the
row determinant. Fix the vacuum pair, the base weight and
coefficient map, the pinned orthonormal frame, and the finite
source translate \(s\). The base level and the varying ray level are
\[
 M_0=g_0g_1,\qquad M=M_0d,
\]
where \(d\) is odd,
squarefree, prime to \(M_0\ell\) and to the conductor and
discriminant exclusions of the host. We use the uniformized
cover with its compatible moduli and adelic frames. Its
coefficient field contains the ray field of level \(M\), the
harmonic coefficients under consideration, and the normalized
values of the splitting characters. A cover with these levels and frames is called admissible for the row at the chosen coefficient place.
\end{paramdefinition}

At the fixed place over \(\ell\), the exponent of \(\ell\) in the level is unchanged. The new marking is a basis of prime-to-\(\ell\) torsion. This torsion is finite étale on the good model, and the compatible bases form a finite étale cover there. Pullback to that cover therefore preserves the good-model condition used in the coefficient bound.

\begin{auxlemma}[Integral finite Fourier expansion]\label{lem:cover-vacuum-fourier}
For a level \(M=M_0d\) of Definition~\ref{hyp:ray-cover-integrality},
let \(\Phi\) be an integral level function on the selected half
lattice, periodic modulo \(M\). If \(k\) generators suffice for
the reduction of the tensor lattice, the vacuum \(v\) and its Heisenberg translates \(W(y)v\) satisfy
\[
 M^k f_\Phi\in\sum_y O\,W(y)v,
\]
with \(O\) the coefficient valuation ring.
The clearing cost at \(\ell\) is at most \(k v_\ell(M_0)\).
\end{auxlemma}
\begin{proof}
Let \(\Lambda\) be the adelic tensor lattice of the vacuum
and write \(\Lambda_e\) and \(\Lambda_{\bar e}\) for its
two projected half lattices. Here \(e\) is the sum of the
idempotents at the places of the level selected by the finite
test vector. A level function \(\Phi\) is supported on
\(\Lambda_{\bar e}\) and is periodic modulo
\(M\Lambda_{\bar e}\). Its values belong to the chosen
valuation ring. Fix a list of \(k\) generators for the
reduction of \(\Lambda\) modulo any positive integer, using
generators for the fixed source and target lattices.
Project these to \(\Lambda_e\) and divide them by \(M\), giving
\(y_i\). The character of \(t\) paired with
\[
 y(\alpha)=\sum_i\alpha_i y_i,
\]
for
\(\alpha\in(\mathbb Z/M\mathbb Z)^k\), is a product of
\(M\)-th roots of unity.

By vacuum self-duality the two half lattices pair perfectly.
For \(t\in\Lambda_{\bar e}\), put \(z_i=\psi(\langle t,y_i\rangle)\).
Then \(z_i^M=1\) and
\[
 \sum_{\alpha}\psi(\langle t,y(\alpha)\rangle)
 =\prod_{i=1}^{k}\sum_{a=0}^{M-1}z_i^a
 =M^k\mathbf 1_{t\in M\Lambda_{\bar e}} .
\]
Indeed, a factor vanishes unless \(z_i=1\); all the \(z_i\) equal one
precisely when \(t\) pairs trivially with the generators. This argument
allows repetitions among the generators. Let \(W(y)\) be
Heisenberg translation by \(y\). Fourier inversion
writes the finite theta vector \(f_\Phi\) as
\[
 M^k f_\Phi
   =\sum_{\alpha}
       \left(\sum_{t\in\Lambda_{\bar e}/M\Lambda_{\bar e}}
             \Phi(t)\,\psi(-\langle t,y(\alpha)\rangle)\right)
             W(y(\alpha))f_0 ,
\]
The coefficients are integral, including at a prime dividing
\(M\). Since \(d\) is prime to \(\ell\), the clearing
exponent in this step is \(k v_\ell(M_0)\).
\end{proof}

\paragraph{Moving to the geometric half lattice.}
The moduli level selects a half lattice at each split
place. Switch the halves that differ from this choice by
the central finite Weil operators. A switch exchanges a
place with its conjugate in the idempotent and carries
each translated vacuum to a translated vacuum in that half,
times the splitting-character scalar. At primes of \(d\)
these switches have identity component at \(\ell\). The
scalar there is integral at \(\ell\) because the local norm factors
are powers of rational residue cardinalities prime to \(\ell\),
and the character uses the norm-normalized values in the
cover field. At primes of \(M_0\) there are finitely many
subsets of places to switch. If \(a_J\) denotes the scalar for such a subset \(J\), a common nonnegative clearing exponent satisfies
\[
 c_{\rm flip}\ge\max_J\{-v_\ell(a_J),0\},
\]
and the associated source
translations form a fixed finite set \(Z\).

For each translate \(sz\), \(z\in Z\), use the split
class representatives away from \(2\ell M\). Concretely
the finite adelic point, after the auxiliary switches, is
the product of a rational isometry \(\gamma\), a finite
point \(h\) supported at split primes away from \(2\ell M\),
and a compact element fixing both the standard lattice and
the level translates of the vacuum. Such a factorization
comes from split-place approximation in the open
stabilizer of the lattice and level vectors. The archimedean
factor is the unique correction to the finite adelic
identity. Write the finite factorization as \(sz=\gamma h k\), with \(k\) the compact factor. At \(\ell\) it gives
\[
 \gamma=(sz)k^{-1}h^{-1},\qquad \gamma^{-1}=hk(sz)^{-1},
\]
with \(h,k\) and their inverses integral there. Thus the
entries of \(\gamma\) and its inverse have a bound on the
\(\ell\)-part of their denominators from \(sz\), even though
the rational isometry varies.

When the Weil kernel is
evaluated at zero, rational automorphy removes \(\gamma\)
and the compact element acts trivially. At infinity this is
a linear substitution in each block of the harmonic
polynomial. It preserves the defining harmonic relations
and weights. In each degree \(h\) its coefficients are
sums of products of \(h\) substitution entries. The finitely
many degrees and the finite set \(Z\) therefore give a
common additional clearing exponent \(c_{\rm sub}\).
The rational source space here is the standard one and
the local entries are tested at all embeddings used in
the cover. These conditions are needed in the split
factorization and the bound on both matrices.

\begin{lemma}[The clearing exponent for the short sections]
\label{lem:theta-short-clearing}
Use the cover and fixed data in Definition~\ref{hyp:ray-cover-integrality}
and the split representative away from \(2\ell M\). At a good frame, the short section for an integral
translate mean and harmonic data of total derivative degree \(h\)
is integral after multiplication by \(\ell^{c_{\rm jet}}\), where
\[
 c_{\rm jet}=v_\ell\!\left(
     M_0^{[E:\mathbb Q]nr+n k_W}(2D_E M_0)^h\right).
\]
Here \(n,r\) are the source and target ranks and \(k_W\) is the
\(\mathbb Z\)-rank of the target marking lattice. The section and
its pinned reading use the same harmonic polynomial and
coefficients of the translate mean.
\end{lemma}

\begin{proof}
Let \(\mathcal T\) be the tuple set of torsion translates of the
source letters at level \(M\). A half-level translate has the
same reading as its representative in \(\mathcal T\): their
difference belongs to the integral lattice of that half, and
their Heisenberg pairing vanishes. The reading of a single
translate at \(t_0\) is the normalized mean with coefficients
\[
 u(t)=|\mathcal T|\,\mathbf 1_{t=t_0}.
\]
In particular its value
in the mean does not acquire a factor of \(|\mathcal T|^{-1}\)
at \(\ell\). Lemma~\ref{lem:cover-vacuum-fourier}, after its indicated
clearing, gives integral linear combinations of these coefficient
arrays.

Let \(R\) be the cover valuation ring and \(P\) the polynomial
in the matrix letters of the moduli frame. The finite module over
\(R\) generated by the coefficients of the array \((u(t)P)_t\)
is torsion-free in the coefficient extension. Choose a basis
\((b_j)_j\) and write the array as \(\sum_j b_j P_j\) with
\(P_j\) arrays over \(R\). Write \(D_\lambda\) for any of the linear operators cutting out
the prescribed harmonic weight over the fraction field of \(R\).
Coordinate extraction in the basis \((b_j)\) gives
\[
 D_\lambda\big((u(t)P)_t\big)
        =\sum_j b_jD_\lambda(P_j)=0
        \quad\Longrightarrow\quad D_\lambda(P_j)=0\quad\text{for all \(j\)}.
\]
Each \(P_j\) has the prescribed weight and derivative degree. By linearity of the jet construction \(J\),
\[
 J\big((u(t)P)_t\big)=\sum_j b_j J(P_j).
\]

The short-section construction on the split representative lattice
supplies the algebraic section and its evaluation by that Weil
kernel average, with the fixed Lie, component, and splitting-character
factors. The hypotheses are the full level marking, the standard
source lattice at \(2\ell M\) with self-duality at \(\ell\), the
integral torsion indicators and constant section, and the prescribed
row degrees. Cuspidality uses the dominant weight and strictly
positive last exponent. In the scalar on this class, the determinant
idele has an ideal representative away from \(2\ell M\).
The normalized unramified character times the square root of its
norm is a local unit at the test place. The remaining inverse
norms and lattice indices, including the index from the split
lattice sandwich, are prime to \(\ell\).

Let
\[
 N=Mu_0
\]
clear the lattice displacement, where \((u_0,2\ell)=1\).
The level index divides \(N^{[E:\mathbb Q]nr}\) and the translate
normalization divides \(M^{n k_W}\) by the standard-lattice and
marking conditions. The latter denominator is harmless for the
indicator array computed above, and may still be included in a
common bound. The good-germ calculation of
Lemma~\ref{lem:regular-theta-jets} costs \((2D_E)^h\).
The additional cost of pullback by \([N]\) is \(N^h\) since
\[
 d[N]=N\,\mathrm{id}
\]
and every derivative monomial has total
degree \(h\). Thus the possible denominator divides
\[
      N^{[E:\mathbb Q]nr} M^{n k_W}(2D_E N)^h.
\]
Its \(\ell\)-valuation is \(c_{\rm jet}\), since \(M/M_0\) and
\(N/M\) are prime to \(\ell\). Its prime-to-\(\ell\) part is a
unit at the frame. This proves the bound there. Compatibility
with degenerating integral frames is proved separately below.
\end{proof}

\paragraph{Good-reduction germs and the connection.}

\begin{lemma}[Denominators of derivatives at a good frame]\label{lem:regular-theta-jets}
For an integral regular germ in the specified theta
trivialization at a good frame, let \(h\) be the order of a jet
in integral tangent directions. If \(a_h\) is its value in the modified polarization gauge and \(R\) is the valuation ring of the frame, then
\[
 (2D_E)^h a_h\in R.
\] The bound descends from a finite extension carrying
the required torsion points.
\end{lemma}

\begin{proof}
Work over a complete valuation ring \(R\) for a good
frame. Use the compatible abelian scheme and the rigidified
symmetric theta bundle. A finite extension carries all
\(N^2\)-torsion used in the short sum and the chart with the
pulled-back frame. Multiplication by \(M\) kills the
level sections. Thus if \([N]t\) is a level section and
\(M\mid N\), then \([N^2]t=0\).
The degree-one section is the integral generator specified
by the frame. The cubical isomorphism
\[
 [N]^*\mathcal L\simeq\mathcal L^{N^2}
\]
preserves the
rigidification and extends over \(R\). It is the normalized
pullback in the short-section formula. After this pullback
the two moves are symmetric theta group translates at
\(N^2\)-torsion. The oddness of \(N\) allows the symmetric
lifts to be chosen of finite order. Each torsion section
extends to the abelian scheme by properness.

The extended torsion section lies in the kernel of the polarization of \(\mathcal L^{N^2}\)
over \(R\), because this equality holds on the generic
fiber and the dual abelian scheme is separated. The
theta group over this section is a multiplicative-group
torsor over the local base, so it has an integral lift.
Write
\[
 g_{\rm fin}=c g_R,
\]
where \(g_{\rm fin}\) is the
finite-order generic lift and \(g_R\) is the integral lift.
Take \(m\) divisible by the order of \(g_{\rm fin}\) and the
order of the underlying torsion section. Then \(g_R^m\) is
a scalar in \(R^\times\), and
\[
  1=g_{\rm fin}^m=c^m g_R^m,\qquad m v(c)=-v(g_R^m)=0 .
\]
Thus \(c\in R^\times\) and both translates extend in the
prescribed normalization. This argument permits
torsion of order divisible by the residue characteristic.
Completing in the integral rigidified trivialization
gives power series over \(R\) for these germs.

\proofstep{Changing the biextension trivialization.}
Write \(B(X,Y)\) for the formal biextension in an integral
chart and \(g(X,Y)\) for a logarithmic gauge over the
fraction field. It vanishes on the axes. The identities
for its two addition defects, transpose defect and CM
balance defect are the logarithms of the corresponding
integral units with constant term one. In particular
their constant first derivatives and their mixed
quadratic coefficients are integral. Write \(\beta\)
for the bilinear coefficient matrix of \(g\), and
\(A(a)\) for the action of \(a\in\mathcal O_E\) on the tangent
space. If the tangent rank is \(m\), integrality of the rigidified mixed coefficients gives
\[
 A(a),\quad \beta^t-\beta,\quad A(a)^t\beta-\beta A(\bar a)\ \in M_m(R).
\]

Choose integral pairs \(x_i,y_i\) giving the separability
element times \(D_E\) in \(E\otimes_{\mathbb Q}E\).
One can take a basis and its trace-dual basis with
the latter multiplied by \(D_E\). They satisfy
\[
 \sum_i x_iy_i=D_E,\qquad \sum_i ax_i\otimes y_i=\sum_i x_i\otimes ay_i\quad(a\in E).
\]
Put
\[
  \gamma=\sum_i A(x_i)^t\beta A(\bar y_i),\qquad
  \beta_0=(\gamma+\gamma^t)/(2D_E).
\]
The latter is symmetric and \(A(a)^t\beta_0=\beta_0 A(\bar a)\).
Put \(T=\gamma-D_E\beta\). The balance defect makes \(T\) integral.
Symmetrizing now gives
\[
 2D_E(\beta-\beta_0)=D_E(\beta-\beta^t)-T-T^t .
\]
Each term on the right is integral in the prescribed frame.

Subtract the symmetric balanced bilinear form \(\beta_0\)
in the logarithmic coordinates of the formal group.
The resulting gauge still satisfies the required identities.
In taking the harmonic value this change contributes
contractions of opposite CM types. These are zero
because the polynomial is in the joint harmonic kernel.
The same contraction rule gives the invariance of the
harmonic value under change of chart with the fixed
tangent and theta frames.

\proofstep{Integral derivatives in the modified gauge.}
Let \(D_u\) be the invariant derivation on the formal group
in an integral tangent vector \(u\), and let \(q(X)\)
be the logarithmic correction obtained by evaluating the
gauge at \((X,[w_0]X)\), with \(w_0\) the fixed integral
polarizing element. Let \(d_0=2D_E\). We have
\(q(0)=D_uq(0)=0\) and
\[
 d_0 D_uD_v q\in R[[X]].
\]

For the last statement, differentiate the two addition
identities at the two identity arguments. The
nonconstant terms are derivatives of the logarithms of
the integral units in the addition defects. Such first
derivatives are integral by \(d\log B=B^{-1}dB\).
The remaining terms are the bilinear coefficient
evaluated on \(u,[w_0]v\) and \(v,[w_0]u\), cleared by \(d_0\).
The algebra generated by \(R[[X]]\) and all
\(d_0D_uq\) for the finitely many tangent directions
has integral constant coefficients. It is stable under
each \(d_0D_u\).

For a regular germ \(f\) and integer
bundle power \(j\), applying \(d_0D_u\) to
\(a\exp(jq)\) acts on the coefficient by
\[
  a\longmapsto d_0D_ua+j(d_0D_uq)a .
\]
Starting from \(a=f\) proves integrality of every
degree-\(h\) derivative after multiplication by \(d_0^h\).
The exponential contributes constant term one, so
no bound on its higher coefficients is used. Apply this
letter by letter and use the integral polynomial
coefficients. A finite extension multiplies normalized
discrete valuations by a positive ramification index,
so the resulting inequality descends to the original
good frame.
\end{proof}

\paragraph{Integral tests at degenerating frames.}\label{par:theta-degenerate-arc}
We use the good-at-\(\ell\) level condition of the
admitted cover and the toroidal moduli model over the
selected coefficient place. The degeneration description
uses the split Raynaud extension and polarized period
lattice after finite extension. The relative toroidal
construction over the complete base turns these data
back into a semiabelian family. \cite{faltings-chai}
For a definite signature
all the test frames have good reduction. In the other
signature consider any integral test frame. By the
valuative criterion on the compactification it extends
to a trait.

If the trait lies in the abelian locus, the
preceding estimate applies. Otherwise, after a finite
extension, split the degeneration data and write its
period pairing as
\(\varpi^{Q(\lambda,\mu)}u_{\lambda,\mu}\) with
unit factors and the integral degeneration form \(Q\).
The period pairing on the arc and its specialization are
\[
 X^{Q(\lambda,\mu)}u_{\lambda,\mu}\ \;\longmapsto\;\ \varpi^{Q(\lambda,\mu)}u_{\lambda,\mu}.
\]
The linear
period and polarization data use the same rescaling.
The integral frame lifts on this arc with constant
tangent coefficients and the same degree-one
theta generator. This uses the derivative
identification for the semiabelian uniformization
and the integral generator at the base point.
Inverting \(\varpi X\) puts the family in the
generic abelian locus. Compatibility of algebraic
evaluation with the arc and its specialization
expresses the value of a scaled section as
\(H(X)/(\varpi X)^b\), where \(H\in R[[X]]\)
and \(b\ge0\).

\begin{lemma}[Integrality after specialization of the toroidal arc]\label{lem:theta-arc-integrality}
On the framed toroidal arc of Paragraph~\ref{par:theta-degenerate-arc},
specializing at \(X=\varpi\), let a scaled section be
\(H(X)/(\varpi X)^b\), with \(H\in R[[X]]\) and \(b\ge0\).
Suppose the section is bounded on the normalized complex frame
and integral at the algebraic good-point tests with the same
clearing power. On the arc and the original trait, respectively, one has
\[
 H\in \varpi^b X^b R[[X]],\qquad H(\varpi)/\varpi^{2b}\in R.
\]
\end{lemma}

\begin{proof}
\proofstep{Test the order at \(X=0\) on the complex frame.} For the complex test,
embed the coefficient fraction field abstractly into
\(\mathbb C\) over the specified coefficient embedding.
Use the holomorphic uniformization of the arc with
\(X=\exp(2\pi i z)\) near the puncture. To compare
the two evaluations, take an affine open chart of the
generic moduli scheme containing the Laurent-series point
and use its algebraic frame. The chart-coordinate
comparison for the holomorphic uniformization of the
split degeneration gives convergent Laurent series for
the arc's chart coordinates. The framed comparison uses
the constant tangent periods and the nonzero degree-one
theta generator of this same arc.

The frame comparison gives convergent Laurent series for the entries of the change-of-frame
matrices and for the ratio of theta generators, with
values equal to the corresponding analytic quantities.
The determinants of the frame matrices and the theta
ratio are invertible formally. Each such denominator germ can be written
\[
 d(X)=X^a u(X),\qquad u(0)\ne0,\qquad d(X)^{-1}=X^{-a}u(X)^{-1}.
\]
After shrinking the punctured disk the last germ converges. In the chart, evaluation
of a regular section is evaluation in the chart frame
followed by its algebraic weight factor. The latter uses
just the entries of the frame change, inverse
determinants and inverse theta ratio. Sums and products
of the convergent readings therefore give the analytic
value with Laurent coefficients equal to the transported
formal coefficients.  \label{cor:theta-cover-integrality-1}

To bound its value, normalize the complex
frame by a Borel change to unit Gram frame.
The cusp section is bounded there in its invariant
norm. Along the arc the diagonal tangent ratios
and the generator ratio in the weight character are
bounded as follows. In the symplectic normalization of
the arc the imaginary period matrix is
\[
 A+\operatorname{Im}(z)Q.
\]

Positivity for all sufficiently large \(\operatorname{Im}(z)\)
implies \(Q\ge0\). For a fixed \(z_0\) in that region the
matrix dominates its positive value at \(z_0\). More explicitly, writing
\[
 Y(z)=A+\operatorname{Im}(z)Q,
\]
the order inequalities for \(\operatorname{Im}z\ge\operatorname{Im}z_0\)
are
\[
 Y(z)\succeq Y(z_0)>0,
 \qquad Y(z)^{-1}\preceq Y(z_0)^{-1},
 \qquad \det Y(z)\ge\det Y(z_0).
\]
The inverse bound controls the tangent ratios and the determinant
bound controls the generator norm in the same period frame. The squared Hermitian length of each constant
tangent vector is computed using the inverse imaginary
period matrix. The squared diagonal ratios from
Gram orthogonalization are bounded by these lengths
times the fixed block scalars.

For the theta generator \(e\), the Gaussian integral over a symplectic period parallelepiped is
\[
 \|e\|^2=\frac{|\kappa|^2}{\sqrt{\det(2\operatorname{Im}\tau)}} ,
\]
where \(\kappa\) is the arc's constant generator scalar. The unit frame uses a generator of squared norm
one. Thus the generator ratio is bounded as well.
The dominant weight uses these ratios only with
nonnegative exponents. Thus the complex Laurent
function is bounded as \(\operatorname{Im}z\to\infty\).
Its negative coefficients vanish by Cauchy's formula.
Injectivity of the coefficient embedding gives
\(X^b\mid H\) in \(R[[X]]\).

\proofstep{The nonarchimedean valuation.}
Holomorphic testing on complex frames bounds the order of vanishing. To use this information in division, we still need the valuation of the coefficient in the nonarchimedean frame.

Write
\[
 H=X^b H_1.
\]
For the second test use the
valuation at \((\varpi)\) in \(R[[X]]\), extended to the
completion of its fraction field. The coefficient formula for this valuation is
\[
 v(\varpi)=1,\qquad v(X)=0,\qquad
 v\!\left(\sum_i a_iX^i\right)=\min_i v(a_i)\quad(a_i\in R).
\]
The arc is now in the abelian locus. The frame still
gives an integral point there, since its tangent and
theta generators lift on the arc. The good-point
bound extends to this valuation by the smooth
integral flag charts. Concretely, on an open
abelian chart the section is a rational function
in integral chart coordinates. If it had a pole
at an integral test with arbitrary residue extension,
one can retain its finite pole order on a point with
algebraic coefficients by smooth lifting of the
finite congruences in that chart. The condition to
stay in the open locus is one of these congruences.
Such a pole contradicts the algebraic good-point tests.

Since \(X\) is a unit in this valuation, the valuation of the
section is \(\min_i v([X^i]H_1)-b\). The good-point tests give
\[
 \min_i v([X^i]H_1)\ge b.
\]
Thus \(H=\varpi^bX^b H_2\)
with \(H_2\in R[[X]]\). The value on the original framed trait
is \(H_2(\varpi)\in R\). Both tests used the section after its
fixed multiplication by a power of \(\ell\).
\label{cor:theta-cover-integrality-2}
\end{proof}

Lemma~\ref{lem:theta-arc-integrality} thus extends the good-point
bound to the degenerating frames of the same cover.

\begin{proposition}[One phase and one denominator on the admitted covers]
\label{prop:theta-uniform-cover-clearing}
Fix the data at the start of this subsection and a finitely generated
integral lattice of input polynomials. There is a phase \(\eta\) from the
base frames and a clearing power \(\ell^c\) such that the comparison
with the algebraic reading of each admitted coefficient is integral
at every prescribed place above \(\ell\), after scaling by
\(\eta\ell^c\). One can take
\[
    c=k v_\ell(M_0)+c_{\rm flip}+c_{\rm sub}+c_{\rm jet}+c_P,
\]
with a common maximum over the degrees of the finitely many
polynomial generators when needed. The varying cover, level function,
twist, places and class label do not enter this choice.
\end{proposition}

\begin{proof}
The source frame in the norm comparison changes by a unitary map
preserving the vacuum lattice. On each component compare the
vacuum identities on a nonzero theta value. If \(c_1,c_2\)
are the two component constants, and \(A\) is the product of
that value, the Gaussian and the nonzero character factor, the
comparison of phases is the comparison of \(c_1 A\) and \(c_2 A\).
The common nonzero factor cancels in the phase ratio:
\[
 \frac{c_1 A}{c_2 A}=\frac{c_1}{c_2}.
\]
This reads the base vacuum phase. The dominant pins
on the cover read these same constants after pullback. Multiplying
its inverse by the inverse phase of the fixed lowest weight vector
in the archimedean harmonic map gives \(\eta\).

The matrix coefficients of that harmonic map on a fixed set of
lattice generators form a finite set of algebraic numbers in each
admitted coefficient field. For each coefficient \(a\), take a
positive integer \(d_a\) for which \(d_a a\) is integral. Then
\[
 c_P=\max_a v_\ell(d_a)
\]
works at all places above \(\ell\),
since \(d_a/\ell^{v_\ell(d_a)}\) is a unit there. The short-section
construction over the field gives the rational section. The local
construction of Lemma~\ref{lem:theta-short-clearing}, extended
by Lemma~\ref{lem:theta-arc-integrality}, gives integral comparison
sections. They have the same pinned adelic reading after the indicated
scaling, so injectivity identifies each with the rational section
at that place.

The determinant twist is a finite-order character on the
norm-one determinant quotient, trivial on the rational points
and the cover congruence subgroup. The indicated ray field gives
its values on the components. Each such value is a root of unity,
so multiplication by it preserves the integral section lattice.
It costs no clearing power. Summing the exponents from Fourier
inversion, the half switch, the change of variables, the short
section and the fixed harmonic map gives the displayed value of
\(c\). \label{cor:theta-cover-integrality-3}
\end{proof}

\subsubsection*{Descent to the original row}
\label{subsec:theta-cover-descent}

The preceding coefficient bounds were obtained in frames on the covers. To measure a coordinate in the original cusp lattice we must also account for the different of the coefficient extension in which those frames were evaluated.

\begin{proposition}[Uniform descent from ray covers]\label{prop:ray-cover-descent}
Suppose a section on the prescribed ray cover descends as a
holomorphic form and is integral at all good frames above
\(\mathfrak q\mid\mathfrak p\mid\ell\) of its coefficient
extension \(L/K_0\). Each coordinate \(c\) in an integral cusp basis at \(\mathfrak p\) satisfies
\[
 c\in\mathfrak D_{L_{\mathfrak q}/(K_0)_{\mathfrak p}}^{-1}
\]
under the chosen local embedding.
Unramified descent preserves integrality. For the universal row
covers contained in the fraction field of the fixed coefficient
DVR \(O\), multiplication by \(\ell^{e_O^2}\) suffices at the
chosen place, where \(e_O=v_O(\ell)\).
\end{proposition}

\begin{proof}
\proofstep{Descent of the holomorphic span.}
Let a form on the cover be invariant under the base level group.
The character relating its adelic and holomorphic realizations is invariant
under that group by the level condition on the splitting
character. Hence the weighted holomorphic form factors
through the base quotient. The bounded cuspidal holomorphic
space of the base is spanned by the complex realizations of its algebraic
cuspidal sections in the prescribed dominant weight with
positive last exponent.

In the application its degree is
at least \(8fr+4\), where \(r\) is the target rank and
\(f\) is the base degree parameter.
Pull the sections back to the cover. Linear independence persists under
extension to \(\mathbb C\). Among the cover pullbacks,
choose a basis over the cover field \(L\). If \(s_L\) lay outside the \(L\)-span of independent pullbacks \(F_1,\ldots,F_t\), then
\[
 \dim_L\langle F_1,\ldots,F_t,s_L\rangle=t+1.
\]
Their complexifications would remain independent, contradicting complex descent. Thus the rational
descent expresses the section in the \(L\)-span of the
original base sections, over the base field \(K_0\).

For the integral statement used on the ray coefficients,
let \(\mathfrak q\mid\ell\) be the selected place of \(L\)
and \(\mathfrak p\) its restriction to \(K_0\).
The base cusp sections integral at \(\mathfrak p\) form a
full lattice over \({\cal O}_{K_0,\mathfrak p}\). The
lattice is defined by the valuation bounds at all good
algebraic frames. It is finite over the localized base ring
by the unit ball bound on the finite-dimensional space.
For any section \(v\) in the base-field space, a common scalar clears its finitely many coordinates, so
\[
 \exists\,a\in\mathcal O_{K_0,\mathfrak p}\setminus\{0\}:\quad av\ \text{lies in the integral section lattice}.
\]
Thus one can use a
free lattice basis \(F_1,\ldots,F_t\). Write the descended
cover section as \(\sum_i y_i F_i\) with \(y_i\in L\),
using the pullbacks for this equality.

\proofstep{Testing the dual lattice by local traces.}
Write
\[
 s_L=\sum_i y_iF_i.
\]
For a base frame \(x\),
pullback commutes with evaluating each \(F_i\). At a place
\(\mathfrak q\) where \(s_L\) is integral, this is true
in every local embedding of \(L_\mathfrak q\) into a common
valued field. At the other places over \(\mathfrak p\), choose
\(b\in\mathcal O_L\) with \(v_{\mathfrak q'}(b)\) large
enough to clear the finitely many \(y_i\). Then every local
evaluation of \(b\,s_L\) on the lift of \(x\) is integral.
The trace section is
\[
    \operatorname{Tr}_{L/K_0}(b\,s_L)
       =\sum_i\operatorname{Tr}_{L/K_0}(by_i)F_i .
\]
Its evaluation at \(x\) is the sum over the embeddings
of \(b\,s_L\), so the base lattice test applies to it.
Since the \(F_i\) form an integral basis, each coefficient
of this trace section is integral at \(\mathfrak p\).

By weak approximation take \(b\)'s approaching any fixed
\(b_\mathfrak q\in\mathcal O_{L_\mathfrak q}\) at \(\mathfrak q\)
and approaching zero at the other local factors. Local
trace and the closedness of the integral base ring give
\[
   \operatorname{Tr}_{L_\mathfrak q/K_{0,\mathfrak p}}
             (b_\mathfrak q y_i)\in
                   \mathcal O_{K_{0,\mathfrak p}}^{\wedge}
      \quad(b_\mathfrak q\in\mathcal O_{L_\mathfrak q}).
\]
This is precisely
\[
 y_i\in\mathfrak D_{\mathfrak q/\mathfrak p}^{-1},
\]
with \(\mathfrak D\) the local different. The test used
all frames at the chosen place, as needed for the full trace.
For an unramified extension the different is the unit ideal.

For a finite prescribed set of unramified places, take one such
local decomposition \(s_L=\sum_i y_{v,i}F_{v,i}\) per place \(v\).
Rescale each \(F_{v,i}\) by a base scalar that is a unit
at the place where it was selected and clears its poles at
the other base places, and rescale \(y_{v,i}\) inversely.
The Chinese remainder theorem gives \(\rho_v\in L\) with
\[
 \sum_v\rho_v=1,
\]
\(\rho_v\) a unit at its place and
divisible by sufficient powers at all the others. Indeed one
may choose all but one by their local congruences and set
the last to one minus their sum.

At a prescribed place \(w\), the choice of powers is
\[
 v_w(\rho_v)\ \ge\
      \max_i\max\{0,-v_w(y_{v,i})\}\quad (v\ne w).
\]
The \(F_{v,i}\) are integral there by the initial rescaling.
For \(v=w\), both \(F_{w,i}\) and \(y_{w,i}\) are integral by
local descent and the triviality of the unramified different.
It follows term by term that the section
\[
 \sum_{v,i}\rho_v y_{v,i}F_{v,i}=s_L
\]
is integral at \(w\).

\proofstep{A uniform bound on the different.}
The descent formula expresses the new denominator through a local different. It remains to bound that different uniformly as the allowed ray levels change.

In the application to the universal row coefficients the
cover can have ramification at \(\ell\). The actual cover
field embeds in the fixed fraction field of the row DVR
\(O\). Take \(e_O=v_{\!O}(\ell)\) with \(v_{\!O}(O^\times)=0\)
and a uniformizer of value one. The restriction of \(v_{\!O}\)
to \(L\) is a positive integral multiple of the discrete
valuation at \(\mathfrak q\). Writing \(a\ge1\) for the integer multiplying the valuation at \(\mathfrak q\), evaluation on \(\ell\) gives
\[
 e_O=a\,v_{\mathfrak q}(\ell)\ge v_{\mathfrak q}(\ell).
\]
A bound of
\(e_O^2\) suffices for the relative different exponent.

For instance, over the maximal unramified subextension
of \(L_{\mathfrak q}/\mathbb Q_\ell\), a uniformizer has
an Eisenstein polynomial of degree \(a\), the absolute
ramification index. The different has the valuation of
its derivative. The nonzero terms of the derivative
have distinct valuation residues modulo \(a\), so its
valuation is at most that of its leading term,
\[
 a-1+a v_\ell(a)<a^2,
\]
including \(a=1\).
Transitivity of the different then bounds the relative
exponent by the same number. Multiplying by
 \(\ell^{e_O^2}\) makes the inverse-different coordinates
integral. In fact, writing \(e_{\rm abs}=v_{\mathfrak q}(\ell)\)
and \(d_{\rm rel}=v_{\mathfrak q}(\mathfrak D_{\mathfrak q/\mathfrak p})\),
\[
 v_{\mathfrak q}(\ell^{e_O^2}y_i)
       \ge e_O^2 e_{\rm abs}-d_{\rm rel}\ge0
\]
by \(d_{\rm rel}<e_{\rm abs}^2\) and \(e_{\rm abs}\le e_O\). These two numerical inequalities give the uniform estimate
\[
 d_{\rm rel}<e_{\rm abs}^2\le e_O^2.
\]
Its role is important even
when the additional level is prime to \(\ell\), since
the coefficient field also contains finite-order values.

\proofstep{The fields used before specialization.}
To see which covers this argument serves, let \(Q\) and
\(m_{\rm ray}\) be the auxiliary places and the depth
for the universal coefficient. The construction takes
the split rational primes below the selected conjugate
half of \(Q\), with the designated prime of the row twist.
The avoidance conditions make their product a squarefree
auxiliary level, prime to \(\ell\), the base level and the
host exclusions.

Write \(K_{\rm ray}\) for the ray class field at the joint level, \(A\) for the finite set of character values of the finite type and distinguished twists, and \(K_{\rm row}\) for the row coefficient field under the fixed embedding. The ray-field condition and coefficient choice give
\[
 K_{\rm ray}\subseteq K_{\rm ray}(A)\subseteq K_{\rm row}.
\]
The uniformized, dominantly
pinned cover is formed over the resulting number
field. Thus both the field conditions for integrality
on the cover and the ramification bound for descent
hold for these choices. In particular this construction
uses the cover field before specializing a character of
the variable depth.

The finite universal data on this cover are indicators
times finite-character values. Let \(A_Q\) be the class-label
group, \(b(g)=\operatorname{detchar}_Q(g)\) the determinant
label and \(S_{a'}\) the section for label \(a'\). Then the
assembled section is
\[
              \Xi_a(g)=\sum_{b\in A_Q}
                     \mathbf1_{b(g)=b}\,S_{a-b}(g).
\]
Each determinant fibre is a union of components by the
ray-field comparison. The indicators preserve the section
lattice, and the finite-character values are integral. Thus
the sum uses the same integrality bound without a denominator
\(|A_Q|\). The prescribed finite type and the trivial
upper-left reductions make \(\Xi_a\) invariant under the
original level group, so the trace descent applies.

The resulting base sections are still bounded at their
chosen base place. Under the faithful row realization,
their coordinate vectors lie in the localization of the
supremum lattice at that place. For a coefficient vector \(x\), localization gives \(x=y/s\) with \(y\) a bounded global vector and \(s\) a base scalar of valuation zero. Under the chosen embedding,
\[
 v_\ell(s)=0\quad\Longrightarrow\quad s^{-1}\in O^\times.
\]
Thus passing from \(y\) to \(x\) costs no further denominator. Taking the additional clearing power to be
\(e_O^2\), the descent yields a bounded row for every
coefficient, with the complex phase and the fixed
harmonic comparison scalar supplied by the pinned
coefficient map. The clearance depends on the datum
and on the fixed \(O\), independently of the level,
depth and class label. The passage from this complex
reading to the actual rational row lattice is given next.
\label{cor:theta-cover-integrality-4}
\end{proof}
